\clearpage
\SourceLeaf{390}
\fancyhead[L]{\small VERSCHIEDENE AUFSÄTZE.}
\pdfbookmark[0]{Verschiedene Aufsätze.}{verschiedene-aufsaetze}
\pdfbookmark[1]{Bestimmung der grössten Ellipse. 1810.}{groesste-ellipse-1810}
\begin{center}
{\Large VERSCHIEDENE AUFSÄTZE.}\par\bigskip
\rule{.85\textwidth}{.3pt}\par
v.\ \textsc{Zach}. Monatliche Correspondenz für Erd- und Himmelskunde. 1810 August.\par
\rule{.85\textwidth}{.3pt}\par\bigskip
BESTIMMUNG DER GRÖSSTEN ELLIPSE\par\medskip
WELCHE DIE VIER SEITEN EINES GEGEBNEN VIERECKS BERÜHRT.\par\bigskip
\rule{.15\textwidth}{.3pt}
\end{center}
\hspace*{\parindent}Die Lage aller Punkte in der Ebene, in welcher das Viereck liegt, bestimme ich durch Abscissen und Ordinaten, indem ich vorerst die Abscissen-Linie und den Anfangspunkt der Abscissen ganz nach Willkür annehme. Das Viereck bestimme ich nicht durch die Winkelpunkte, sondern durch die Punkte, wo jenes Seiten von den aus dem Anfangspunkte der Abscissen auf diese gefällten Perpendikeln geschnitten werden. Diese Perpendikel seien \(a,a',a'',a'''\), und ihre Neigungen gegen die Abscissen-Linie \(A,A',A'',A'''\), folglich die Coordinaten der erwähnten vier Durchschnittspunkte
\[
\begin{array}{ll}
a\cos A,&a\sin A\\
a'\cos A',&a'\sin A'\\
a''\cos A'',&a''\sin A''\\
a'''\cos A''',&a'''\sin A'''
\end{array}
\]
